\setcounter{footnote}{76}

\printedpage{73}
% head: A. Grothendieck : Dérivateurs, Ch. XII, Édition des Universités de Montpellier II et Paris VII
% foot: Version 7 décembre 2002

\msp{page 105}

I will again take $(X,U,\{t\})$ such that $t$ is a final object of $X$, but if we want $q\notin W$, we cannot suppose $t'$ a final element of $X'$. Instead of this, I am going to suppose

\begin{enumerate}[label=\alph*)]
\item $U$, $U'$ (and also $X$, having final object $t$) $W$-aspheric, hence $i\in W$ (and even $i$ $W$-aspheric), $p\in W$, but
\item $X'$ not $W$-aspheric, hence $i',q\notin W$.
\end{enumerate}

For this, let us note that
\[ U'/t'\simeq U'/p_!(e_U)\overset{\text{def}}{=}\underline{\pi}_0(U/U'). \]
The corollary to proposition 7 (cf. p. 85), not yet proved, tells us that $i'\in W$ (i.e. here, $X'$ $W$-aspheric) if and only if $i'$ is $W$-aspheric, i.e. $U'/t'$ is $W$-aspheric. We shall admit it here. Everything thus comes down to finding an arrow
\[ p:U\longrightarrow U' \]
in Cat with the following properties.

\begin{enumerate}[label=\alph*)]
\item $U$ and $U'$ are $W$-aspheric,
\item $\underline{\pi}_0(U/U')$ is not $W$-aspheric.
\end{enumerate}

We shall take up the preceding example again
\[ f:X\longrightarrow Y \]
(p. 101-103), taking
\[ U=X,\qquad \underline{\pi}_0(U/U')=Y, \]

\msp{page 106}

so it all comes down, starting from the given example $X\xrightarrow{f}Y$, to constructing an arrow
\[ \pi:Y\longrightarrow U' \tag*{$(*)$} \]
in Cat which satisfies the conditions

\begin{enumerate}[label=\alph*)]
\item $\pi$ is a local isomorphism,
\item $U'$ is $W$-aspheric.
\end{enumerate}

We shall then take
\[ p=\pi f. \]
But from the way $X$ has been constructed from a 1-connected and not $\infty$-connected $Y$, otherwise arbitrary, we are finally reduced to finding $(*)$ satisfying a), b), and moreover

\begin{enumerate}[label=\alph*), start=3]
\item \printedpage{74}
% head: A. Grothendieck : Dérivateurs, Ch. XII, Édition des Universités de Montpellier II et Paris VII
% foot: Version 7 décembre 2002
$Y$ is 1-connected, but not $\infty$-connected.
\end{enumerate}

(NB We now suppose $W=W_\infty$.) It amounts to the same thing to find

% printed with a vertical rule at the left of this block
\begin{quote}
$U'$ an \textit{aspheric} object of Cat, $F\in\operatorname{Ob}U'^\wedge$ such that $U'/F$ is 1-connected and not $\infty$-connected.
\end{quote}

One thinks for example of simply taking an open subcategory $U'_0$ of $U'$, which is 1-connected and not $\infty$-connected. Let us start with such a category $U'_0$ - for example a homotopy realization in Cat, of the sphere $S^2$. We immerse it as an open in an aspheric $U'$, by adding a final object to it. We have won.

\msp{page 107}

Let us make a statement in due form.

\textbf{Proposition 11.} \textit{One can find a situation} $(*)$ (p. 124 \texttt{[rather 104]}) \textit{with}

\begin{enumerate}[label=\textit{\alph*)}]
\item \textit{$t$ a final object of $X$.}
\item \textit{$U$, $U'$ (and $X$) $W_\infty$-aspheric, hence $i,p\in W_\infty$ (and even $i$ $W_\infty$-aspheric).}
\item \textit{$X'$ not $W_\infty$-aspheric, hence $i',q\notin W_\infty$.}
\end{enumerate}

For this it suffices to start with a 1-connected object $Y$ of Cat, to adjoin to it a final object $u'_0$, which gives $U'$, to take $U=\underline{\mathrm{Ch}}(Y;y_0,-)$, $p:U\longrightarrow U'$ the composite $U\xrightarrow{\text{target}}Y\underset{\text{can.}}{\overset{\text{incl.}}{\hra}}U'$, and $X$ deduced from $U$ by adding a final object $t$ to it. We then have $U'/t'\simeq Y$, whence
\[ X'\ W_\infty\text{-aspheric (i.e. }i'\in W_\infty)\Longleftrightarrow Y\ W_\infty\text{-aspheric.} \]
We have the desired counterexample by taking $Y$ (1-connected but) not $\infty$-connected.

\subsection*{2.4\quad Application to open Dwyerian immersions.}

Let us give all the same the \textit{positive} statement we were after. % [sic] l'enoncé -> l'énoncé, missing accent

\textbf{Theorem 3.} \textit{Let $(X,U,T)$ be a triple $\in\mathcal{T}$ ($U$ open in $X$ in} Cat\textit{, $T=X\backslash U$), and for every arrow $p:U\longrightarrow U'$ an arrow of} Cat\textit{, the cocartesian square in} Cat

\[
\begin{tikzcd}
U \arrow[r, hook, "i"] \arrow[d, "p"'] & X \arrow[d, "q"'] & T \arrow[l, dotted, hook'] \arrow[d, dotted, "\text{iso}"] \\
U' \arrow[r, hook, "i'"] & X' & T'. \arrow[l, dotted, hook']
\end{tikzcd}
\]

\printedpage{75}
% head: A. Grothendieck : Dérivateurs, Ch. XII, Édition des Universités de Montpellier II et Paris VII
% foot: Version 7 décembre 2002

\msp{page 108}

\textit{Suppose that for every $t\in\dot{T}=\dot{U}$, the presheaf of sets}
\[ i^*(t)=(u\longmapsto\operatorname{Hom}(u,t))\in U^\wedge \]
\textit{is totally $W$-aspheric (def. 1, p. 88)} (\footnote{\textbf{NB} One says under these conditions that $i_0:U\longrightarrow\bar{U}$ is \textit{strictly $W$-aspheric} (cf. def. 2 next page).}) \textit{- e.g. that the category $U/t$ is totally $W$-aspheric (prop. 8, page 93). (NB the case where the $U/t$ have a final object, cf. prop. 9 a), i.e. the one where the inclusion $U\longrightarrow\bar{U}$ admits a right adjoint, is the one envisaged by Dwyer, that of open Dwyer immersions} (\footnote{I know of no other case where th. 3 applies - for lack of knowing of a totally $W$-aspheric category which does not have a final object, nor of an object $F$ of a category $U^\wedge$ which is totally $W$-aspheric without being representable, i.e. $U/F$ has a final object. Now, I know that \textit{filtering} categories are totally $W$-aspheric if $W$ satisfies the axiom of direct limits.})\textit{.) Then for every $p:U\longrightarrow U'$, the open immersion}
\[ i'_0:U'\longrightarrow\overline{U'} \]
\textit{is $W$-aspheric (and even, just like $i_0$, strictly $W$-aspheric), and a fortiori it is} \texttt{[in]} \textit{$W$. Consequently (prop. 6} \texttt{[p. 81]}\textit{) we have the following.}

\[
\begin{aligned}
&\text{\textcircled{a}}\qquad && i\in W && \Longleftrightarrow\quad && i'\in W\quad && \textit{assuming } \mathrm{W}(7\text{ bis}).\\
&\text{\textcircled{b}}\qquad && p\in W && \Longleftrightarrow\quad && q\in W\quad && \textit{assuming } \mathrm{W}(7).\ (\footnotemark)
\end{aligned}
\]
\footnotetext{Proof page 112.}

\textbf{Corollary 1.} \textit{Assume} W(7)\textit{. If $i:U\hra X$ is an open immersion such that $i_0:U\hra\bar{U}$ is totally} \texttt{[strictly suffices]} \textit{$W$-aspheric, then $i:U\longrightarrow X$ is a $W$-cofibration (``ultrastrong''), by which I mean this : for every arrow $U'\xrightarrow{f}U''$ in $U\backslash$}Cat\textit{,}

\msp{page 109}

\textit{such that $f\in W$, the arrow $g:X'\longrightarrow X''$ which is deduced from it by cobase change from the object $X$ of $U\backslash$}Cat\textit{, is also in $W$.}

\textsc{Proof of the corollary} : we apply theorem 3 to the open immersion $U'\overset{i'}{\hra}X'$ deduced from $U\overset{i}{\hra}X$ by the cobase change $U\xrightarrow{f}U'$, noting in addition that the hypothesis of theorem 3 on $i$ is preserved in passing from $i$ to $i'$.

We must all the same introduce adequate terminology :

\textbf{Definition 2 :}

\begin{enumerate}[label=\alph*)]
\item An arrow $f:X\longrightarrow Y$ is said to be \textit{strictly $W$-aspheric} if for every $y\in\operatorname{Ob}Y$, $f^*(y)=(x\longmapsto\operatorname{Hom}_Y(f(x),y))$ is a totally $W$-aspheric object of $X^\wedge$ (cf. def. 1, page 88). It is said to be \textit{totally $W$-aspheric} if for every $y\in\operatorname{Ob}Y$, $X/y$ ($\simeq X/f^*(y)$) is a totally $W$-aspheric object of Cat.

\item \printedpage{76}
% head: A. Grothendieck : Dérivateurs, Ch. XII, Édition des Universités de Montpellier II et Paris VII
% foot: Version 7 décembre 2002
An open immersion $i:U\longrightarrow X$ is said to be \textit{$W$-Dwyerian} (resp. \textit{totally} \texttt{[$W$-]} \textit{Dwyerian}) if the induced immersion $i_0:U\longrightarrow\bar{U}$ is strictly (resp. totally) $W$-aspheric. A closed immersion $j:T\longrightarrow X$ is said to be \textit{$W$-Dwyerian} (resp. \textit{totally $W$-Dwyerian}) if and only if $j^\circ:T^\circ\longrightarrow X^\circ$ is an open $W$-Dwyerian immersion (resp totally $W$-Dwyerian).

\item An arrow $i:U\longrightarrow X$ of Cat (not necessarily an immersion)

\msp{page 110}

is called a \textit{$W$-cofibration,} if for every arrow $f:U'\longrightarrow U''$ in $U\backslash$Cat, the arrow $g:X'\longrightarrow X''$ which is deduced from it by cobase change by $X$ (object of $U\backslash$Cat) is in $W$ when $f$ is.
\end{enumerate}

\textbf{NB} This notion of $W$-cofibration is stable under composition and under cobase change, and the isomorphisms of Cat are $W$-cofibrations. If $f$ is a $W$-cofibration, $f^\circ$ is one and conversely.

(The corollary is stated by saying that an open or closed immersion which is $W$-Dwyerian is a $W$-cofibration.)

\begin{enumerate}[label=\alph*), start=4]
\item An arrow $i:U\longrightarrow X$ in Cat is called a \textit{trivial $W$-cofibration}, or a \textit{couniversal $W$-equivalence,} if for every $U'$ in $U\backslash$Cat, i.e. every arrow $p:U\longrightarrow U'$ in Cat, the arrow deduced from $i$,
\[ i':U'\longrightarrow X'=U'\amalg_U X'\quad\texttt{[rather }U'\amalg_U X\texttt{]}\ \ , \]
is in $W$. (NB It is immediate that this indeed implies that $i$ is a $W$-cofibration.)
\end{enumerate}

(This notion of trivial $W$-cofibration is also stable under composition and under cobase change, and the isomorphisms in Cat are such.)

\msp{page 111}

\textbf{Corollary 2.} (Assuming W(7).) \textit{Let $i:U\longrightarrow X$ be an open $W$-Dwyerian immersion. For $i$ to be a trivial $W$-cofibration, it is necessary and sufficient that $i\in W$, or what amounts to the same thing} (cf. page 82), \textit{that $\dot{T}\hra T$ be in $W$.}

This is obviously necessary, and sufficiency follows from theorem 3 a).

\textbf{Open questions.}

\begin{enumerate}[label=\alph*)]
\item I do not know whether an open immersion $i$ which is a $W$-cofibration, or even a trivial $W$-cofibration, is necessarily $W$-Dwyerian (\footnote{and perhaps even a Dwyer immersion?}). But by the corollary to proposition 7 (p. 85), I know that this is the case if we suppose $T$ isomorphic to the final category $e$. At least this follows : For not only $i:U\longrightarrow X$, but also all the open immersions $i_1:U\longrightarrow X_1$ which are deduced from it by the procedure of prop. 7 b) by means of arbitrary arrows $T_1\longrightarrow T$ with target $T$ in Cat, to be $W$-cofibrations, it is necessary and sufficient that $i$ be Dwyerian. % [sic] toute au moins -> tout au moins, agreement error

\item \printedpage{77}
% head: A. Grothendieck : Dérivateurs, Ch. XII, Édition des Universités de Montpellier II et Paris VII
% foot: Version 7 décembre 2002
Is there a good notion of a Dwyerian immersion, applicable to all immersions

\msp{page 112}

and including the case of open immersions and closed immersions?
\end{enumerate}

\textsc{Proof} of theorem 3 (p. 107). It suffices to prove that if $i$ is $W$-Dwyerian (i.e. $i_0$ strictly $W$-aspheric). then $i'$ is so as well - for it will follow that $i'$ is $W$-aspheric and a fortiori $\in W$. We must therefore verify that for every $t'\in\dot{U}'=\dot{T}$, the element $i^*(t')=(u'\longmapsto\operatorname{Hom}_{X'}(u',t'))$ of $U'^\wedge$ is totally $W$-aspheric. Now $t'$ comes from a unique element $t$ of $\dot{T}$, and we know that we have % [sic] asphérique). alors -> asphérique), alors, period before lower-case word

\[ i'^*(t)=p_!(i^*(t))\quad\texttt{[rather }i'^*(t')=p_!(i^*(t))\texttt{]},\qquad\text{where }\ p_!:U^\wedge\longrightarrow U'^\wedge. \]

By hypothesis, $F=i^*(t)$ is a totally $W$-aspheric element of $U^\wedge$, which means that for every arrow $p:U\longrightarrow U'$ in Cat with source $U$, $p_!(F)$ is a $W$-aspheric element of $U'$. In fact, we have the immediate lemma (by transitivity $(p'p)_!\simeq p'_!p_!$). % [sic] souce -> source, missing letter

\textbf{Lemma.} \textit{Let $U$ be in} Cat\textit{, and $F$ in $U^\wedge$. If $F$ is totally $W$-aspheric (resp. if $U/F$ is totally $W$-aspheric), then for every arrow $p:U\longrightarrow U'$ in} Cat \textit{with source $U$, $F'=p_!(F)$ is a totally $W$-aspheric object of $U'^\wedge$ (resp.}

\msp{page 113}

\textit{$U'/F'$ is a totally $W$-aspheric object of} Cat\textit{.)}

\textbf{Corollary 3} (to theorem 3). \textit{If $i:U\longrightarrow X$ is an open or closed $W$-Dwyerian immersion (resp. totally $W$-Dwyerian), then for every arrow $p:U\longrightarrow U'$ in} Cat\textit{, $i':U'\longrightarrow X'=U'\amalg_U X'$} \texttt{[rather $U'\amalg_U X$]} \textit{is so too.}

(NB The non-resp. case is essentially theorem 3. It does not depend on any additional axiom on $W$, but is of interest only in view of the equivalences \textcircled{a} and \textcircled{b} (p. 108) which result from it, and which are valid assuming W(7 bis) resp. W(7) only.)

\textbf{Comment on motivation :} open or closed $W$-Dwyerian immersions, or totally $W$-Dwyerian ones, are stable under cobase change (by the preceding corollary), but do not at all seem to be so under composition. In this respect, ``Dwyer immersions'' (where one supposes the existence of the right adjoint to $i_0:U\longrightarrow\bar{U}$, in the case of an open immersion) behave better, being stable under composition. (\textbf{NB} It does not seem

\msp{page 114}

clear either that strictly $W$-aspheric arrows, or totally $W$-aspheric ones, are stable under composition - sorry!) My motivation for wearing myself out over the $W$-Dwyerian notion is twofold.

\begin{enumerate}[label=\alph*)]
\item I would like to have conditions at once comprehensible, and as general as possible, on an immersion (open e.g.), for it to be a $W$-cofibration, a visibly important property (just like the dual notion of $W$-fibration). % [sic] toute comme -> tout comme, agreement error
\end{enumerate}